\section{Requisitos do Sistema de Detecção de Anomalias}
\label{sec:requisitos_ml}

Esta seção especifica o detector de anomalias. Ele parte do RSSI que o
\textit{gateway} mede em cada quadro, um dos metadados de recepção descritos na
Seção~\ref{subsec:camada_fisica_lora} e registrados pela infraestrutura da
Seção~\ref{sec:requisitos_infra}, e complementa a autenticação da
Seção~\ref{sec:requisitos_seguranca}, que não indica de onde o sinal partiu. O
modelo é estimado e avaliado com os dados do \textit{dataset} Tour Perret.

%%%%%%%%%%%%%%%%%%%%%%%%%%%%%%%%%%%%%%%%%%%%%%%%%%%%

\subsection{Requisitos funcionais do modelo de detecção}

\begin{enumerate}

    \item[RF-M1.] O sistema deve prever o RSSI de cada quadro um passo à frente,
    a partir dos quadros anteriores do mesmo nó sensor no mesmo
    \textit{gateway}, com um modelo por canal de frequência.

    \item[RF-M2.] O sistema deve prever o RSSI com RMSE menor que o do
    \textit{random walk} por canal.

    \item[RF-M3.] O sistema deve sinalizar como anômalo o quadro cujo erro de
    previsão passa do limiar do seu canal, calibrado no conjunto de validação.

\end{enumerate}

%%%%%%%%%%%%%%%%%%%%%%%%%%%%%%%%%%%%%%%%%%%%%%%%%%%%

\subsection{Requisitos não funcionais do modelo de detecção}

\begin{enumerate}

    \item[RNF-M1.] O modelo deve atingir o desempenho do RF-M2 com no máximo 14
    dias de quadros de treino.

    \item[RNF-M2.] O modelo deve ter no máximo dois parâmetros estimados por
    canal, para caber num dispositivo de borda.

    \item[RNF-M3.] Todo resultado deve ser reproduzível a partir do
    \textit{dataset} bruto, conferido por resumo SHA-256.

\end{enumerate}

\subsection{Rastreabilidade dos requisitos de detecção}

O Quadro~\ref{qua:rastreabilidade_ml} liga cada requisito à seção do
Capítulo~\ref{cap:desenvolvimento} em que ele é verificado.

\begin{quadro}[htb]
	\centering
	\caption{Rastreabilidade dos requisitos do modelo de detecção}
	\label{qua:rastreabilidade_ml}
	\begin{tabular}{lll}
		\hline
		\textbf{Requisito} & \textbf{Onde é verificado} & \textbf{Situação} \\
		\hline
		RF-M1 & Seção~\ref{subsec:ml-res-frequencia} & atendido \\
		RF-M2 & Seção~\ref{subsec:ml-res-frequencia} & atendido no RMSE e no MAE \\
		RF-M3 & Seção~\ref{subsec:ml-res-deteccao} & atendido com desvios sintéticos \\
		RNF-M1 & Seção~\ref{subsec:ml-res-treino} & atendido na validação \\
		RNF-M2 & Seção~\ref{subsec:ml-res-frequencia} & atendido \\
		RNF-M3 & Seção~\ref{subsec:ml-res-serie} & atendido \\
		\hline
	\end{tabular}
	\legend{Fonte: elaborado pelos autores.}
\end{quadro}
